% 颜色
% 介绍混色：% 80% StructureColor + 20% black

% === LiuYin 主题配色 ===
\ifdefstring{\Ysy@ycolor}{LiuYin}{
	\definecolor{StructureColor}{HTML}{97c6c0} % #97c6c0（清茶绿）
	\definecolor{MainColor}{HTML}{4df8e8} % #4df8e8（死星过载）
	\definecolor{EmphColor}{HTML}{e26e1b} % #e26e1b（火焰橙）
	\definecolor{LinkColor}{HTML}{475d7b} % #475d7b（塞纳河畔的春水）
}{\relax}

% === Catppuccin Latte 主题配色 ===
\ifdefstring{\Ysy@ycolor}{CatppuccinLatte}{
	\definecolor{StructureColor}{HTML}{7287fd}  % lavender：柔和的薰衣草蓝
	\definecolor{MainColor}{HTML}{04a5e5}       % sky：清亮的天蓝色
	\definecolor{EmphColor}{HTML}{d20f39}       % red：深酒红色
	\definecolor{LinkColor}{HTML}{209fb5}       % sapphire：沉静理性的蓝青色
}{\relax}

% === SPA 主题配色 -> 替换为冷自然的配色方案 ===
\ifdefstring{\Ysy@ycolor}{SPA}{
	\definecolor{StructureColor}{HTML}{403990}  % 深蓝
	\definecolor{MainColor}{HTML}{80A6E2}       % 浅蓝
	\definecolor{EmphColor}{HTML}{F46F43}       % 橘红
	\definecolor{LinkColor}{HTML}{FBDD85}       % 浅黄
}{\relax}

% === ElegantBlue 主题配色（来自ElegantBook） ===
\ifdefstring{\Ysy@ycolor}{Elegant}{
	\definecolor{StructureColor}{RGB}{60,113,183} % 经典蓝色
	\definecolor{MainColor}{RGB}{0,166,82} % 活力绿色
	\definecolor{EmphColor}{RGB}{255,134,24} % 明亮橙色
	\definecolor{LinkColor}{RGB}{0,174,247} % 亮蓝色
}{\relax}

% === 纯黑论文学术配色 ===
\ifdefstring{\Ysy@ycolor}{OhMyPaper}{
	\definecolor{StructureColor}{HTML}{000000}  
	\definecolor{MainColor}{HTML}{000000}       
	\definecolor{EmphColor}{HTML}{000000} 
	\definecolor{LinkColor}{HTML}{000000}
}{\relax}

% === 蓝色星空 ===
\ifdefstring{\Ysy@ycolor}{StarFieldBlue}{
	\definecolor{StructureColor}{HTML}{252a50}  % 信标蓝
	\definecolor{MainColor}{HTML}{3b5b8a}    % 已见蓝   
	\definecolor{EmphColor}{HTML}{f05627}  % 橙红
	\definecolor{LinkColor}{HTML}{d01c1f} % 热焰红
}{\relax}

% === 粉蓝天空 ===
\ifdefstring{\Ysy@ycolor}{SkyPinkBlue}{
	\definecolor{StructureColor}{HTML}{555d8b}  % 反射蓝
	\definecolor{MainColor}{HTML}{766c9b}    % 石斛之花   
	\definecolor{EmphColor}{HTML}{eca4b8}  % 胭脂扣
	\definecolor{LinkColor}{HTML}{f7b9c2} % 眉开眼笑
}{\relax}

% === 绯色天空 ===
\ifdefstring{\Ysy@ycolor}{SkyScarlet}{
	\definecolor{StructureColor}{HTML}{d16d7c}  % 沙漠玫瑰
	\definecolor{MainColor}{HTML}{f7786b}    % 桃红回响
	\definecolor{EmphColor}{HTML}{f39477}  % 活力亮橙   
	\definecolor{LinkColor}{HTML}{8c9cc1} % 薰衣草光
}{\relax}

% === 漫光天空 ===
\ifdefstring{\Ysy@ycolor}{SkyRelax}{
	\definecolor{StructureColor}{HTML}{4c55bc}  % 惬意逸蓝
	\definecolor{MainColor}{HTML}{ad69a2}    % 鹿蛇草   
	\definecolor{EmphColor}{HTML}{fc8f9b}  % 海螺壳
	\definecolor{LinkColor}{HTML}{ffbe98} % 绒毛桃子
}{\relax}

% === Mondrian ===
\ifdefstring{\Ysy@ycolor}{Mondrian}{
	\definecolor{StructureColor}{HTML}{0000ff}  % 蓝
	\definecolor{MainColor}{HTML}{25d4d0}    % 浅蓝   
	\definecolor{EmphColor}{HTML}{f20d0d}  % 红
	\definecolor{LinkColor}{HTML}{ffff00} % 黄
}{\relax}

% === 碧屿浮光 ===
\ifdefstring{\Ysy@ycolor}{EmeraldGleam}{
	\definecolor{StructureColor}{HTML}{094c54}  % 深绿
	\definecolor{MainColor}{HTML}{349d89}    % 墨绿  
	\definecolor{EmphColor}{HTML}{38bfa7}  % 蓝绿
	\definecolor{LinkColor}{HTML}{d8e6c4} % 黄绿
}{\relax}

% === 融入寂静中 ===
\ifdefstring{\Ysy@ycolor}{FadeintoSilence}{
	\definecolor{StructureColor}{HTML}{4c3d3e}  % 棕黑 Dusty Cocoa
	\definecolor{MainColor}{HTML}{383b4e}    % 蓝黑  
	\definecolor{EmphColor}{HTML}{798e9d}  % 蓝灰 Slate Blue
	\definecolor{LinkColor}{HTML}{bab2a7} % 灰黄 Soft Greige
}{\relax}

% === 六月的雨的模样（紫色主题） ===
\ifdefstring{\Ysy@ycolor}{JuneRainfall}{
	\definecolor{StructureColor}{HTML}{5f5baf}  % Misty Violet
	\definecolor{MainColor}{HTML}{9a64c5}    % 
	\definecolor{EmphColor}{HTML}{568de4}  % Rainy Blue  
	\definecolor{LinkColor}{HTML}{ba9fe1} % Rain Purple
}{\relax}